\section{Ricorsione e Algoritmi (Lab 11)}

\subsection{La Ricorsione: Funzioni che chiamano se stesse}
Una funzione è ricorsiva quando "usa se stessa" per risolvere un problema.
Il segreto è avere sempre un **Caso Base** (dove mi fermo) e un **Passo Ricorsivo** (dove vado avanti).

\subsubsection{Esempio: Numeri Triangolari}
La somma dei numeri da 1 a N. `triangle(3) = 3 + 2 + 1 = 6`.

\begin{lstlisting}[language=JavaScript]
function triangle(n) {
    if (n == 1) return 1;          // Caso base: mi fermo a 1
    return n + triangle(n - 1);    // Passo ricorsivo: n + il resto
}

console.log(triangle(6)); // 21
\end{lstlisting}

\subsection{Lo Stack (La Pila dei Piatti)}
\begin{concept}{Visualizzazione Stack Tracing}
Ogni volta che chiami una funzione, il computer aggiunge un "piatto" (Record) sulla pila della memoria.
\begin{enumerate}
    \item Chiamo `triangle(3)` -> Stack: [3]
    \item Lui chiama `triangle(2)` -> Stack: [3, 2]
    \item Lui chiama `triangle(1)` -> Stack: [3, 2, 1] -> Ritorna 1 (POP!)
    \item `triangle(2)` riceve 1, fa 2+1=3 -> Ritorna 3 (POP!)
    \item `triangle(3)` riceve 3, fa 3+3=6 -> Ritorna 6 (POP!)
\end{enumerate}
\end{concept}

\subsection{Tail Recursion (Ricorsione in Coda)}
\begin{concept}{Ottimizzazione Coda}
Per non riempire la memoria (Stack Overflow), passo il risultato parziale come argomento.
In questo modo il computer riutilizza lo stesso spazio.
\end{concept}

\begin{jsbox}
function triangleCoda(n, s=0) {
    if (n < 1) return s;           // Ritorno l'accumulatore
    return triangleCoda(n - 1, s + n); // Passo tutto al prossimo
}
\end{jsbox}

\subsection{Esercizi Classici}

\subsubsection{1. Cifre di un Numero}
\begin{concept}{Logica: Modulo e Divisione}
Come scomporre un numero (es. 123)?
\begin{itemize}
    \item **Ultima cifra:** `n % 10`
    \item **Tolgo l'ultima:** `Math.floor(n / 10)`
\end{itemize}
\end{concept}

\textbf{A. Stampa Cifre (stampaCifreRec)}
\begin{jsbox}
function stampaCifreRec(n) {
    if (n < 10) {
        console.log(n); // Caso base
        return;
    }
    stampaCifreRec(Math.floor(n / 10)); // Stampa le precedenti
    console.log(n % 10); // Stampa l'ultima
}
\end{jsbox}

\textbf{B. Array di Cifre (arrayCifreRec)}
\begin{lstlisting}[language=JavaScript]
function arrayCifreRec(n) {
    if (n < 10) return [n];
    
    let resto = arrayCifreRec(Math.floor(n / 10));
    resto.push(n % 10);
    return resto;
}
\end{lstlisting}

\subsubsection{2. Fibonacci (N-esimo)}
Calcolare solo un numero specifico della sequenza: $F_n = F_{n-1} + F_{n-2}$.
\begin{lstlisting}[language=JavaScript]
function fibonacci(n) {
    if (n <= 1) return n; // 0 o 1
    return fibonacci(n - 1) + fibonacci(n - 2);
}
\end{lstlisting}

\subsubsection{3. Fibonacci (Array)}
Restituire l'array dei primi N numeri di Fibonacci.
\begin{lstlisting}[language=JavaScript]
function fibonacciPrimiN(n) {
    if (n == 0) return [];
    if (n == 1) return [0];
    if (n == 2) return [0, 1];
    
    // Ottengo i precedenti
    let prec = fibonacciPrimiN(n - 1);
    // Calcolo il nuovo: somma degli ultimi due
    let nuovo = prec[prec.length - 1] + prec[prec.length - 2];
    prec.push(nuovo);
    return prec;
}
\end{lstlisting}

\subsubsection{2. Ricerca Binaria (Su array ordinato)}
Invece di cercare uno per uno, dividiamo sempre a metà.
\begin{lstlisting}[language=JavaScript]
function binarySearch(arr, el, sx, dx) {
    if (sx > dx) return -1; // Non trovato
    
    let centro = Math.floor((sx + dx) / 2);
    
    if (arr[centro] == el) return centro; // Trovato!
    
    if (arr[centro] > el) {
        // Cerca nella metà sinistra
        return binarySearch(arr, el, sx, centro - 1);
    } else {
        // Cerca nella metà destra
        return binarySearch(arr, el, centro + 1, dx);
    }
}
\end{lstlisting}
